% =============================================================================
% Pre-registration v1 - The Causal Safety Audit (Paper 2)
% Deposited on OSF. This file is the source of the deposited PDF; the PDF and
% its SHA-256 in this directory are the authoritative record.
% =============================================================================
\documentclass[11pt,a4paper]{article}

\usepackage[margin=2.4cm]{geometry}
\usepackage[T1]{fontenc}
\usepackage{lmodern}
\usepackage{amsmath,amssymb}
\usepackage{booktabs}
\usepackage{enumitem}
\usepackage[hidelinks]{hyperref}
\usepackage{parskip}

\title{\textbf{Pre-registration v1}\\[0.3em]
The Causal Safety Audit: Counterfactual Stress-Testing for\\
Deep-Learning MRI Reconstruction}

\author{Dat Tat Mai\thanks{School of Science, Engineering \& Technology, RMIT
University Vietnam, Ho Chi Minh City, Vietnam. ORCID
0009-0006-5256-5646.}
\and Thai Viet Pham\thanks{School of Computing Technologies, RMIT University,
Melbourne, VIC, Australia. ORCID 0009-0003-3707-5081.}
\and Thi Nguyen\thanks{VinFuture.}}

\date{2026-09-08}

\begin{document}
\maketitle

\begin{center}
\fbox{\parbox{0.92\textwidth}{\small
\textbf{Status at deposit.} No reconstruction has been run. No pilot has been
run. No experimental result of any kind exists at the time of this deposit.
The thresholds in Section~\ref{sec:thresholds} were fixed from convention and
prior literature, with no contact with reconstruction output. The code as of
this deposit is archived at Zenodo concept DOI
\href{https://doi.org/10.5281/zenodo.22648224}{10.5281/zenodo.22648224}
(release \texttt{v0.1.0}, version DOI 10.5281/zenodo.22648225).
}}
\end{center}

% -----------------------------------------------------------------------------
\section{Study rationale and question}
% -----------------------------------------------------------------------------

Deep-learning MRI reconstruction is evaluated almost universally with global
image scores (PSNR, SSIM) and with data-consistency residuals. Both are
insensitive, by construction, to content the measurement did not constrain: a
reconstruction may remove a small lesion and still report an excellent global
score and a residual at the noise floor. This study asks whether that failure
mode occurs at rates of clinical relevance, and whether the standard metrics
detect it.

The audit inserts a lesion on the \emph{measurement} side of the true
multi-coil forward operator, $Y_{\mathrm{cf}} = Y + \mathbf{A}\ell$, so that
the factual and counterfactual arms share the measured noise realisation,
sampling mask and coil sensitivities exactly. Both arms are reconstructed with
the same model and the same seed. Every lesion-attributable quantity is a
within-pair difference.

% -----------------------------------------------------------------------------
\section{Hypotheses}
% -----------------------------------------------------------------------------

\begin{description}[leftmargin=3.2em,style=nextline]
\item[\textbf{H1}] The silent erasure rate for lesions of
$\le 10\,\mathrm{mm}^3$ at $R \ge 6$ exceeds the pre-specified rate of
clinical relevance $\pi_0$ in at least one model family.
\emph{Directional; one-sided.}

\item[\textbf{H2}] Global PSNR and SSIM do not separate erased from preserved
pairs: the standardised mean difference in each score lies within the
equivalence margin, \emph{and} the AUC for classifying erasure from the score
alone lies within the margin about $0.5$. \emph{Equivalence hypothesis; both
conditions required.}

\item[\textbf{H3}] ROI contrast and per-lesion detectability $z$ do separate
erased from preserved pairs, with classification AUC above the pre-specified
floor. \emph{Directional.}

\item[\textbf{H4}] The false-structure count in the lesion-absent arm
increases with acceleration $R$, and is highest for the generative family.
\emph{Directional; trend and family contrast.}
\end{description}

H2 is a null result by construction. It is therefore tested as an equivalence
hypothesis against margins fixed in advance. A merely non-significant
difference will be reported as inconclusive and \emph{not} as support for H2.

% -----------------------------------------------------------------------------
\section{Design}
% -----------------------------------------------------------------------------

\paragraph{Data.} The multi-coil brain subset of fastMRI, which provides fully
sampled raw $k$-space with coil sensitivities recoverable from the calibration
region. Sequences are T1-weighted, T2-weighted and FLAIR axial slices at
1.5\,T and 3\,T. Retrospective undersampling of already-processed magnitude
images is excluded as a validity requirement. Lesion annotations come from
fastMRI+; slices used for insertion are those with no annotation within the
audited region.

\paragraph{Synthetic lesion bank.} Spherical or ellipsoidal perturbations with
volumes $1$, $5$, $10$ and $27\,\mathrm{mm}^3$; three contrast levels spanning
the fastMRI+ interquartile range for the matching sequence; both signs of
contrast where the sequence admits them (hyperintense on FLAIR and T2,
hypointense on T1). Sites are drawn from white matter, deep grey matter and
the cortical ribbon, avoiding ventricles and skull.

\paragraph{Models under audit.} Three families, audited \emph{as released}:
(i) End-to-End Variational Network (physics-based unrolled, learned data
consistency); (ii) an image-domain complex U-Net baseline; (iii) a score-based
diffusion reconstructor with a data-consistency step, evaluated with common
random numbers across the two arms of each pair. No model is retrained on the
lesion bank, and no lesion-bearing slice enters any training or calibration
step.

\paragraph{Conditions.} $R \in \{4, 6, 8\}$ with equispaced Cartesian masks;
random masks as a secondary condition.

\paragraph{Unit of analysis.} The lesion pair. Pairs are nested in slices and
slices in subjects, and one host slice carries several lesions, so pairs are
not independent.

% -----------------------------------------------------------------------------
\section{Pre-specified thresholds}
\label{sec:thresholds}
% -----------------------------------------------------------------------------

These values are fixed by this deposit. They are normative choices drawn from
convention and prior literature, not estimates from data. They are recorded
machine-readably in \texttt{configs/thresholds.yaml} in the archived
repository.

\begin{center}
\begin{tabular}{llp{8.4cm}}
\toprule
\textbf{Symbol} & \textbf{Value} & \textbf{Meaning and justification} \\
\midrule
$z_{\mathrm{det}}$ & $3.0$ &
Reference detectability required for a lesion to enter the erasure analysis.
The channelised Hotelling observer per-lesion $z$ is a normal deviate; $3.0$
is the conventional ``clearly visible'' anchor. The Rose criterion ($5.0$) was
considered and rejected: it would exclude most of the $\le 10\,\mathrm{mm}^3$
bank that H1 concerns, and so would answer a different question. \\
\addlinespace
$z_{\mathrm{miss}}$ & $2.0$ &
Reconstruction detectability below which the lesion counts as missed. Set
strictly below $z_{\mathrm{det}}$ by design: pairs falling in the
indeterminate band $[2.0, 3.0)$ are counted as \emph{neither} erased nor
preserved, so borderline pairs cannot inflate the erasure rate. \\
\addlinespace
$z_{\mathrm{fs}}$ & $4.0$ &
False-structure threshold, in standard deviations of the null-space error over
background tissue. Candidates are screened over many connected components per
slice; $4$\,sd controls the per-image false-positive count. A candidate must
additionally have a volume of at least the smallest lesion in the bank. \\
\addlinespace
$\pi_0$ & $0.05$ &
Silent erasure rate taken as the floor of clinical relevance for sub-$10\,
\mathrm{mm}^3$ lesions at $R \ge 6$. A domain judgement, not an estimate. \\
\bottomrule
\end{tabular}
\end{center}

\paragraph{Definition of erasure.} A pair is \emph{erased} when
\begin{equation}
z_i^{\mathrm{ref}} \ge z_{\mathrm{det}}
\quad\text{and}\quad
z_i^{\mathrm{recon}} < z_{\mathrm{miss}} .
\end{equation}
A pair is \emph{silently erased} when, in addition, the global PSNR of the
lesion-present reconstruction lies within the central $95\%$ interval of PSNR
over lesion-absent reconstructions of the same slice set under the same model
and $R$, so that the reported score gives no indication that anything is
missing.

\paragraph{Other fixed analysis constants.}
Equivalence margin on the standardised mean difference, $\pm 0.2$;
AUC equivalence interval for H2, $[0.45, 0.55]$;
AUC floor for H3, $0.80$;
family-wise $\alpha = 0.05$ with Holm multiplicity correction;
one-sided $\alpha = 0.025$ for H1;
PSNR interval level $0.95$;
regularisation rule $\lambda = \sigma_\eta^2 / \mathrm{var}(|X|)$ with a
sensitivity analysis across one decade either side;
target half-width of the $95\%$ interval for a rate of $\pi_0$, $0.02$.

% -----------------------------------------------------------------------------
\section{Statistical analysis plan}
% -----------------------------------------------------------------------------

\paragraph{Primary endpoint.} The silent erasure rate for lesions of
$\le 10\,\mathrm{mm}^3$ at $R \ge 6$, per model.

\paragraph{Dependence.} Every rate is estimated with a confidence interval
from a cluster bootstrap that resamples \emph{subjects}. Every model-level
contrast is fitted as a mixed-effects model with random intercepts for subject
and slice. The paired structure of the two arms is preserved throughout.

\paragraph{H1.} Tested against $\pi_0 = 0.05$ with a one-sided
cluster-bootstrap test at $\alpha = 0.025$, model by model.

\paragraph{H2.} For each of PSNR and SSIM, the standardised mean difference
between erased and preserved pairs is tested with two one-sided tests against
a margin of $\pm 0.2$, using cluster-robust standard errors; equivalence is
declared only if both tests reject. The AUC of the global score as a
classifier of erasure is estimated with a subject-level bootstrap interval and
equivalence to $0.5$ is declared if that interval lies within $[0.45, 0.55]$.
Both conditions are required.

\paragraph{H3.} Compares the AUC of $z$, of ROI contrast and of
$t_{\mathcal{N}}$ as classifiers of erasure against the floor of $0.80$ and
against the AUC of PSNR, with subject-level bootstrap intervals for the paired
difference.

\paragraph{H4.} Fits the false-structure count as a mixed-effects Poisson
model with $R$ as an ordered factor and family as a fixed effect, and tests
the trend in $R$ and the family contrast.

\paragraph{Multiplicity.} H1 to H4 form one family of four primary tests,
controlled with the Holm procedure at a family-wise level of $0.05$. Analyses
by lesion volume, contrast, site and sequence, the sensitivity analyses over
$\lambda$ and observer internal noise, and the Case~B analysis on real
annotated lesions are \emph{secondary} and are reported with intervals, not
tests.

\paragraph{Reporting.} All per-pair readouts are released as a table carrying
subject, slice, site, volume, contrast, model, $R$, mask and seed identifiers,
together with provenance columns (git commit, config hash, checkpoint hash,
data checksums), so that every figure and every rate in the paper can be
recomputed from it.

% -----------------------------------------------------------------------------
\section{Sample size and the design pilot}
\label{sec:samplesize}
% -----------------------------------------------------------------------------

The lesion bank is sized so that the half-width of the $95\%$ interval for a
rate of $\pi_0$ is at most $0.02$ after inflation by the design effect
$1 + (m-1)\rho$, with $m$ the pairs per subject and $\rho$ the intra-subject
correlation.

$\rho$ is not known at the time of this deposit and will be estimated from a
design pilot of 20 subjects. The following constraints are fixed \emph{now}:

\begin{enumerate}[leftmargin=2em]
\item The sizing \emph{rule} above is fixed by this deposit. Only the realised
number of pairs per model-by-$R$ cell remains to be computed from it.
\item The pilot estimates $\rho$ \textbf{and nothing else}. No endpoint, no
erasure rate, no hypothesis test and no AUC will be computed on pilot data.
\item The 20 pilot subjects are \textbf{excluded} from the confirmatory
analysis set.
\item The realised pairs-per-cell will be recorded in a dated addendum to this
registration, posted \emph{before} any confirmatory reconstruction is run.
\end{enumerate}

This sequencing is deliberate. It places the timestamp on the thresholds and
the analysis plan before any contact whatsoever with reconstruction output,
rather than after a pilot has been seen.

% -----------------------------------------------------------------------------
\section{Exclusions, stopping and deviations}
% -----------------------------------------------------------------------------

\paragraph{Slice-level exclusions.} Slices carrying a fastMRI+ annotation
within the audited region are excluded from insertion. Slices failing the
integrity check (calibration region unusable, coil sensitivity estimation
non-convergent) are excluded and counted.

\paragraph{Pair-level exclusions.} Pairs whose reference detectability falls
below $z_{\mathrm{det}}$ do not enter the erasure analysis: the lesion was not
detectable even in the fully sampled reference, so its absence in the
reconstruction is not attributable to the model. These are counted and
reported.

\paragraph{Indeterminate pairs.} Pairs with $z_i^{\mathrm{recon}}$ in
$[z_{\mathrm{miss}}, z_{\mathrm{det}})$ are neither erased nor preserved.
Their count is reported per cell.

\paragraph{No optional stopping.} The number of pairs per cell is fixed before
the confirmatory run and analysis is performed once, after the full run
completes.

\paragraph{Deviations.} Every departure from this document is recorded in
\texttt{CHANGELOG.md} in the archived repository, with its date and reason,
and reported in the paper.

% -----------------------------------------------------------------------------
\section{Outcome-neutral conditions}
% -----------------------------------------------------------------------------

The following must hold for the study to be interpretable, independently of
which hypotheses are supported:

\begin{enumerate}[leftmargin=2em]
\item The forward operator passes its adjoint test.
\item Lesion insertion passes the two-path equality test (insertion through
the operator equals insertion in image space followed by the operator).
\item The filtered projectors satisfy the closed-form and transfer-identity
tests, and $t_{\mathcal{R}} = 1$ under hard data consistency.
\item The channelised Hotelling observer does not exceed the ideal-observer
bound.
\item Both arms of every pair share the seed; seeds are varied across
replicates.
\end{enumerate}

Failure of any of these is a defect to be fixed before the confirmatory run,
not a result.

% -----------------------------------------------------------------------------
\section*{Materials}
% -----------------------------------------------------------------------------

Code: \url{https://github.com/dmai287/csa-audit}, archived at Zenodo concept
DOI \href{https://doi.org/10.5281/zenodo.22648224}{10.5281/zenodo.22648224}.
Data: fastMRI and fastMRI+, under their respective data-use terms. Derived
lesion banks are released only if those terms permit; otherwise the generation
script alone is released.

\end{document}
